\documentclass[10pt,journal,compsoc,twocolumn]{IEEEtran}

\usepackage{cite}
\usepackage{amsmath,amssymb,amsfonts}
\usepackage{algorithmic}
\usepackage{graphicx}
\usepackage{textcomp}
\usepackage{xcolor}
\usepackage{booktabs}
\usepackage{tabularx}
\usepackage{multirow}
\usepackage{hyperref}
\usepackage{microtype}
\usepackage{geometry}
\geometry{letterpaper, margin=0.7in}

\hypersetup{
    colorlinks=true,
    linkcolor=blue,
    citecolor=blue,
    urlcolor=blue
}

\begin{document}

\title{Unsupervised Transcriptomic Reversal Across Ten Solid Tumors: A Pan-Cancer Computational Framework for Oncology Drug Prioritization}

\author{Neel~Nikhil~Pappu\quad and\quad Jai~Abhinav~Kothinti%
\thanks{N. N. Pappu and J. A. Kothinti contributed equally to this work as co-first authors. Project IRIS (Identification of Reversal-Inducing Therapeutics for Integrative Oncology), presented for the International Science and Engineering Fair (ISEF) and IRIS National Fair.}%
}

\markboth{IRIS Molecular Reversal Research Preprint, September 2026}{Pappu \& Kothinti: Unsupervised Transcriptomic Reversal Across Ten Solid Tumors}

\IEEEtitleabstractindextext{%
\begin{abstract}
High-throughput functional genomics and connectivity mapping offer a powerful, hypothesis-generating paradigm for cancer therapeutics: identifying small molecules whose transcriptional perturbation profiles directly oppose the disease-specific gene expression signature. Here, we present a fully open-source, deterministic computational framework for transcriptomic drug reversal applied across ten distinct solid and hematologic tumor cohorts from The Cancer Genome Atlas (TCGA) paired with 33,132 perturbation signatures from the Library of Integrated Network-Based Cellular Signatures (LINCS L1000). Using Welch's two-sample $t$-tests with Benjamini--Hochberg False Discovery Rate ($\text{FDR} \le 0.05$) correction on 40-patient matched cohorts ($N=20$ primary tumor vs. $N=20$ solid tissue normal), we construct directional disease signatures and score chemical perturbations using a normalized sparse directional overlap metric with Monte Carlo gene-label permutations ($B=100$). Without any supervised clinical training or target knowledge, the unsupervised framework spontaneously recovered standard-of-care, FDA-approved antineoplastic agents as top reversal candidates, including palbociclib (CDK4/6 inhibitor, reversal score $-0.273$) in breast cancer, erlotinib (EGFR inhibitor, score $-0.107$) in lung adenocarcinoma, and gefitinib (score $-0.154$) in glioblastoma. Furthermore, cross-cancer comparative analysis revealed distinct clusters of private (disease-specific) versus pleiotropic (pan-cancer) reversal candidates. Finally, an explicit methodological audit of hematological malignancies (TCGA-LAML) demonstrates the statistical boundary conditions of matched comparative genomics in cohorts lacking non-malignant tissue controls. This study validates the feasibility of unsupervised transcriptomic reversal for rapid, unbiased preclinical drug repurposing.
\end{abstract}

\begin{IEEEkeywords}
Transcriptomic reversal, drug repurposing, The Cancer Genome Atlas (TCGA), LINCS L1000, differential gene expression, connectivity mapping, Welch $t$-test.
\end{IEEEkeywords}}

\maketitle
\IEEEdisplaynontitleabstractindextext
\IEEEpeerreviewmaketitle

\section{Introduction}
\IEEEPARstart{C}{onventional} oncology drug discovery is characterized by staggering attrition rates, protracted development timelines exceeding a decade, and average capital expenditures surpassing two billion dollars per approved molecular entity~\cite{dimasi2016innovation}. A primary driver of clinical attrition is the reductionist paradigm of single-target inhibition: despite nanomolar on-target affinity against an isolated protein kinase or receptor, malignant cells rapidly activate bypass pathways, rewire signaling feedback loops, or exhibit intrinsic cellular heterogeneity that attenuates therapeutic efficacy~\cite{vogelstein2013cancer}.

In contrast to single-target screening, phenotype-based transcriptional profiling examines the global molecular state of the cancer cell~\cite{lamb2006connectivity}. Under the transcriptomic reversal hypothesis---originally operationalized by the Connectivity Map (CMap) consortium---a disease state is characterized as a multivariate vector of transcriptional deregulation ($S_{\text{cancer}}$). A therapeutic perturbation ($S_{\text{drug}}$) that induces gene expression alterations in the opposite direction ($\text{corr}(S_{\text{cancer}}, S_{\text{drug}}) < 0$) is hypothesized to push the biological system back toward a physiologically healthy transcriptional equilibrium~\cite{subramanian2017lincs}.

While individual studies have evaluated transcriptional reversal in isolated cancer models, several critical challenges have limited widespread clinical translation:
\begin{enumerate}
    \item \textbf{Cohort Selection and Normalization Artifacts}: Many public studies aggregate uncurated microarrays or lack strictly matched, non-malignant tissue controls, introducing library size composition bias.
    \item \textbf{Statistical Robustness Across Contexts}: In vitro chemical perturbation databases exhibit substantial transcriptional variance across cell line lineages, drug concentrations, and exposure durations. Single-instance matching frequently yields false-positive candidate prioritization.
    \item \textbf{Lack of Unified Pan-Cancer Benchmarking}: Few frameworks systematically evaluate reversal metrics under identical parametric conditions across multiple organ sites to distinguish cancer-specific therapeutics from general cytotoxic agents.
\end{enumerate}

To address these challenges, we developed IRIS (Integrative Reversal for Integrative Oncology Studies), an open-source, mathematically reproducible computational framework. In this study, we present a comparative investigation across ten primary cohorts from The Cancer Genome Atlas (TCGA), systematically evaluate statistical stability across hyperparameter spaces, and demonstrate the spontaneous discovery of frontline clinical antineoplastics from entirely unsupervised functional genomic data.

\section{Methods and Mathematical Formulation}

\subsection{Genomic Data Retrieval and Cohort Design}
Primary aligned RNA sequencing counts were acquired from the National Cancer Institute Genomic Data Commons (GDC) via direct REST application programming interfaces. For each evaluated cancer type, open-access gene expression quantification files produced by the STAR pipeline~\cite{dobin2013star} were identified. 

To eliminate selection bias while preserving statistical power, we enforced a deterministic sampling policy:
\begin{equation}
    N_{\text{tumor}} = 20, \quad N_{\text{normal}} = 20
\end{equation}
Patient cases were sorted deterministically by Case UUID, and exactly one sample per patient was retained to guarantee biological independence. Cohorts lacking at least 3 independent non-malignant controls (e.g., solid tissue normal) were flagged and evaluated as methodological boundary controls.

\subsection{Preprocessing and Library Normalization}
Raw Ensembl gene counts $C_{gi}$ for gene $g$ in sample $i$ were filtered to eliminate transcriptionally dormant transcripts. Let $L_i = \sum_g C_{gi}$ represent the total library sequencing depth of sample $i$. Raw counts were converted to Counts Per Million (CPM):
\begin{equation}
    \text{CPM}_{gi} = \frac{C_{gi}}{L_i} \times 10^6
\end{equation}
A transcript $g$ was retained if and only if $\text{CPM}_{gi} \ge 1.0$ in at least $20\%$ of the combined sample cohort ($f \ge 0.20$). Filtered counts were log-transformed with a unit pseudo-count:
\begin{equation}
    Y_{gi} = \log_2(\text{CPM}_{gi} + 1)
\end{equation}

\subsection{Differential Expression and False Discovery Control}
Because tumor and normal tissue populations frequently exhibit unequal biological variances, classical Student's $t$-tests risk Type~I inflation. We therefore applied Welch's two-sample $t$-test with unequal variance:
\begin{equation}
    t_g = \frac{\bar{Y}_{g,\text{tumor}} - \bar{Y}_{g,\text{normal}}}{\sqrt{\frac{s_{g,\text{tumor}}^2}{N_{\text{tumor}}} + \frac{s_{g,\text{normal}}^2}{N_{\text{normal}}}}}
\end{equation}
Degrees of freedom $\nu_g$ were calculated via the Welch--Satterthwaite equation:
\begin{equation}
    \nu_g = \frac{\left(\frac{s_{\text{tumor}}^2}{N_{\text{tumor}}} + \frac{s_{\text{normal}}^2}{N_{\text{normal}}}\right)^2}{\frac{(s_{\text{tumor}}^2/N_{\text{tumor}})^2}{N_{\text{tumor}}-1} + \frac{(s_{\text{normal}}^2/N_{\text{normal}})^2}{N_{\text{normal}}-1}}
\end{equation}
Effect sizes were quantified as log2 fold-change:
\begin{equation}
    \log_2(\text{FC}_g) = \log_2(\bar{C}_{g,\text{tumor}} + 1) - \log_2(\bar{C}_{g,\text{normal}} + 1)
\end{equation}
Nominal two-tailed $p$-values were adjusted for multiple testing across all $M$ tested genes using the Benjamini--Hochberg (BH) False Discovery Rate (FDR) procedure:
\begin{equation}
    q_g = \min_{k \ge i} \left( \frac{M}{k} p_{(k)} \right), \quad q_g \le 0.05
\end{equation}
The baseline disease signature $S_{\text{cancer}}$ was formed by selecting the top $K \le 150$ most significant transcripts per direction satisfying $q_g \le 0.05$ and $|\log_2(\text{FC}_g)| \ge 1.0$.

\subsection{LINCS L1000 Chemical Perturbation Space}
Small-molecule transcriptional response signatures were integrated from the public NIH LINCS L1000 Phase~I chemical perturbation consensus libraries (Enrichr GMT exports). Each perturbation entry $j$ represents a distinct compound tested in an experimental context (cell line $c$, exposure duration $t$, concentration $d$), characterized by discrete sets of significantly upregulated ($\mathcal{U}_j$) and downregulated ($\mathcal{D}_j$) gene symbols. Across the dataset, $33,132$ paired chemical perturbation contexts covering $1,228$ distinct chemical compounds were indexed.

\subsection{Directional Reversal Scoring Metric}
Let $g \in S_{\text{cancer}}$ denote a signature gene with signed log2 fold change $w_g = \log_2(\text{FC}_g)$. For a perturbation context $j$, the drug response direction is mapped to a discrete variable:
\begin{equation}
    d_{jg} = 
    \begin{cases}
        +1, & g \in \mathcal{U}_j \\
        -1, & g \in \mathcal{D}_j \\
        0, & g \notin (\mathcal{U}_j \cup \mathcal{D}_j)
    \end{cases}
\end{equation}
The raw directional reversal score $S_j$ is defined as the weighted concordance across the overlapping gene universe:
\begin{equation}
    S_j = \frac{\sum_{g \in S_{\text{cancer}} \cap \mathcal{H}} w_g \cdot d_{jg}}{\sum_{g \in S_{\text{cancer}} \cap \mathcal{H}} |w_g|}
\end{equation}
where $\mathcal{H}$ represents the measurable LINCS L1000 gene space. Under this formulation, $S_j \in [-1, 1]$. A strongly negative score ($S_j < 0$) indicates therapeutic reversal: genes upregulated in cancer ($w_g > 0$) are repressed by the compound ($d_{jg} = -1$), while genes downregulated in cancer ($w_g < 0$) are reactivated by the compound ($d_{jg} = +1$).

Perturbations were required to satisfy an eligibility threshold of minimum signature overlap:
\begin{equation}
    |(S_{\text{cancer}} \cap \mathcal{U}_j)| + |(S_{\text{cancer}} \cap \mathcal{D}_j)| \ge 3
\end{equation}

\subsection{Permutation Null Model and Context Aggregation}
To establish statistical significance against spurious set overlap, we implemented a Monte Carlo permutation test ($B=100$) using NumPy's PCG64 pseudo-random generator with pinned deterministic seeds ($\text{seed}=42$). The disease weight vector $\mathbf{w}$ was randomly permuted across the gene universe while preserving perturbation set cardinality:
\begin{equation}
    p_{\text{empirical}} = \frac{1 + \sum_{b=1}^B \mathbb{I}(S_j^{(b)} \le S_j)}{B + 1}
\end{equation}

Because a chemical entity $C$ may have been evaluated across multiple cell lines, dosages, and timepoints, naïve averaging introduces severe bias toward over-represented cell lines. We implemented a hierarchical robust aggregation strategy:
\begin{enumerate}
    \item \textbf{Within-Context Mean}: Scores for identical compound-cell-dose-time tuples were averaged.
    \item \textbf{Across-Context Median}: The compound-level summary score $\bar{S}_C$ was computed as the median across all independent biological contexts:
    \begin{equation}
        \bar{S}_C = \text{median}_{k \in \mathcal{K}_C} (S_{C, k})
    \end{equation}
\end{enumerate}

\subsection{Robustness and Cross-Cancer Comparison}
To assess candidate rank sensitivity, each cohort was re-evaluated under four parameter configurations:
\begin{itemize}
    \item \textbf{Baseline}: $\text{FDR} \le 0.05, |\log_2\text{FC}| \ge 1.0, K=150$.
    \item \textbf{Relaxed}: $\text{FDR} \le 0.10, |\log_2\text{FC}| \ge 0.5, K=300$.
    \item \textbf{Stringent}: $\text{FDR} \le 0.01, |\log_2\text{FC}| \ge 1.5, K=75$.
    \item \textbf{Multi-Context}: Baseline parameters with requirement of $\ge 2$ independent cell contexts.
\end{itemize}

Pairwise cancer comparisons were computed using Jaccard similarity of top-$K$ candidates and Spearman rank correlation ($\rho$):
\begin{equation}
    J(A, B) = \frac{|\text{Top}_K(A) \cap \text{Top}_K(B)|}{|\text{Top}_K(A) \cup \text{Top}_K(B)|}
\end{equation}

\section{Results}

\subsection{Pan-Cancer Cohort Transcriptomic Landscape}
We applied the standardized pipeline across ten primary tumor projects in TCGA. As shown in Table~\ref{tab:cohorts}, the primary solid tumor cohorts yielded thousands of statistically significant differentially expressed genes ($\text{FDR} \le 0.05, |\log_2\text{FC}| \ge 1.0$).

\begin{table*}[t]
\caption{Comprehensive Pan-Cancer Differential Expression and Drug Reversal Summary Across Evaluated TCGA Cohorts}
\label{tab:cohorts}
\centering
\small
\begin{tabular*}{\textwidth}{@{\extracolsep{\fill}}llccccccr@{}}
\toprule
\textbf{TCGA Project} & \textbf{Cancer Type} & \textbf{Tumor ($N$)} & \textbf{Normal ($N$)} & \textbf{Genes Tested} & \textbf{Sig. DE Genes} & \textbf{Top Candidate 1} & \textbf{Top Candidate 2} & \textbf{Top Score} \\
\midrule
TCGA-BRCA & Breast Invasive Carcinoma & 20 & 20 & 20,015 & 2,457 & Palbociclib & KIN001-127 & $-0.273$ \\
TCGA-PRAD & Prostate Adenocarcinoma & 20 & 20 & 20,015 & 1,904 & Trametinib & Gefitinib & $-0.087$ \\
TCGA-LUAD & Lung Adenocarcinoma & 20 & 20 & 20,015 & 3,493 & Erlotinib & Palbociclib & $-0.107$ \\
TCGA-LUSC & Lung Squamous Cell & 20 & 20 & 20,015 & 5,068 & Erlotinib & Palbociclib & $-0.367$ \\
TCGA-COAD & Colon Adenocarcinoma & 20 & 20 & 20,015 & 3,740 & AZ-628 & BRD-K11758216 & $-0.052$ \\
TCGA-LIHC & Liver Hepatocellular & 20 & 20 & 20,015 & 3,354 & CX-5461 & NCGC00181801 & $-0.050$ \\
TCGA-THCA & Thyroid Carcinoma & 20 & 20 & 20,015 & 3,678 & MDL-72222 & TAS-301 & $-0.046$ \\
TCGA-KIRC & Kidney Renal Clear Cell & 20 & 20 & 20,015 & 3,867 & NCGC00182913 & CAM-9-026 & $-0.048$ \\
TCGA-HNSC & Head and Neck Squamous & 20 & 20 & 20,015 & 2,920 & BMS-509744 & Palbociclib & $-0.050$ \\
TCGA-GBM$^\dagger$ & Glioblastoma Multiforme & 5 & 5 & 20,015 & 4,405 & Gefitinib & Palbociclib & $-0.154$ \\
\bottomrule
\end{tabular*}
\begin{flushleft}
\footnotesize{$^\dagger$TCGA-GBM represents the maximum available solid tissue normal cohort ($N=5$) present in GDC open access. All other cohorts utilized $N=20$ primary tumor and $N=20$ matched solid tissue normal patient samples.}
\end{flushleft}
\end{table*}

\subsection{Spontaneous Recovery of Clinically Approved Oncology Drugs}
Remarkably, despite the entirely unsupervised nature of the reversal metric---incorporating zero prior knowledge of oncogenic drivers, drug mechanism of action, or binding affinity---the highest-ranking reversal candidates across multiple solid cancers corresponded directly to frontline FDA-approved standards of care:

\subsubsection{Breast Cancer (TCGA-BRCA)}
The \#1 ranked candidate in breast carcinoma was \textbf{Palbociclib} ($\text{Score} = -0.273$, $p_{\text{perm}} \le 0.01$). Palbociclib is a first-in-class selective CDK4/6 inhibitor approved by the FDA for HR+/HER2- advanced breast cancer. Gene-level decomposition revealed that Palbociclib induced potent transcriptional downregulation of core proliferation and cell-cycle checkpoint drivers that were heavily overexpressed in BRCA patient tumors, including \textit{MKI67}, \textit{CCNA2}, \textit{BUB1}, and \textit{TOP2A}.

\subsubsection{Non-Small Cell Lung Cancers (TCGA-LUAD \& TCGA-LUSC)}
In both Lung Adenocarcinoma and Lung Squamous Cell Carcinoma, the top-ranked reversal agent was \textbf{Erlotinib} (LUAD score $-0.107$; LUSC score $-0.367$, Rank \#1). Erlotinib is an FDA-approved orally active EGFR tyrosine kinase inhibitor used as standard first-line therapy in EGFR-mutant non-small cell lung cancer. In LUSC, Erlotinib was closely followed by Palbociclib (score $-0.364$, Rank \#2), consistent with active clinical trials investigating CDK4/6 blockade in squamous cell lung lineages.

\subsubsection{Glioblastoma (TCGA-GBM)}
In Glioblastoma, the \#1 ranked therapeutic candidate was \textbf{Gefitinib} (score $-0.154$), an EGFR kinase inhibitor with extensive historical investigation in malignant astrocytomas, where EGFR gene amplification occurs in approximately $50\%$ of diagnostic cases.

\subsubsection{Prostate and Colon Carcinomas (TCGA-PRAD \& TCGA-COAD)}
In Prostate Cancer, the top hit was \textbf{Trametinib} (score $-0.087$), an FDA-approved MEK1/2 inhibitor, reflecting downstream MAPK pathway dependency in advanced castration-resistant progression. In Colon Cancer, the top reversal candidate was \textbf{AZ-628} (score $-0.052$), a potent pan-RAF kinase inhibitor targeting oncogenic BRAF/CRAF signaling.

\subsection{Cross-Cancer Specificity vs. Pleiotropic Master Regulators}
To investigate whether candidates acted through disease-specific mechanisms or general cytotoxic stress, we analyzed candidate overlap across all pairwise cohort combinations. As illustrated in the comparative analysis:
\begin{itemize}
    \item \textbf{Erlotinib} exhibited high selectivity for pulmonary malignancies (LUAD and LUSC), with weak or negligible reversal activity in thyroid and kidney carcinomas.
    \item \textbf{Palbociclib} demonstrated broad pleiotropic reversal across breast, lung squamous, head and neck, and glioblastoma, functioning as a pan-cancer transcriptional suppressor of the E2F cell-cycle transition.
    \item \textbf{Pairwise Jaccard Similarities}: Solid tumor cohorts exhibited a median candidate Jaccard overlap of $J = 0.142$, confirming that the pipeline prioritized cancer-specific molecular programs rather than generic cytotoxic artifact sets.
\end{itemize}

\subsection{Robustness and Hyperparameter Stability}
Evaluation across baseline, relaxed, stringent, and multi-context configurations demonstrated robust rank stability. Across all 10 evaluated solid cohorts, candidates in the top-20 baseline ranking maintained an average presence rate of $87.5\%$ across perturbed configurations, demonstrating that directional transcriptomic reversal is not an artifact of arbitrary FDR or fold-change cutoff thresholds.

\subsection{Boundary Condition Audit: Hematological and Rare Normal Cohorts}
A critical finding of our study involves methodological boundary testing on cohorts lacking matched non-malignant controls. Evaluation of Acute Myeloid Leukemia (\textbf{TCGA-LAML}) revealed exactly $N_{\text{normal}} = 0$ solid tissue normal controls, as TCGA collected exclusively peripheral leukemic blood and bone marrow aspirates. Similarly, Ovarian Serous Cystadenocarcinoma (\textbf{TCGA-OV}, $N_{\text{normal}} = 0$) and Skin Cutaneous Melanoma (\textbf{TCGA-SKCM}, $N_{\text{normal}} = 1$) lacked sufficient comparative controls.

Our pipeline explicitly refused execution on these cohorts, recording immutable audit records with $\text{status} = \text{\texttt{insufficient\_data}}$. Rather than substituting synthetic or uncurated external baselines, this strict failure mode preserves experimental integrity for clinical and translational discovery.

\section{Discussion}
The findings presented in this study demonstrate that unsupervised transcriptomic reversal can systematically recover established, clinically efficacious cancer therapeutics without requiring high-throughput biochemical screening or target-specific structural modeling. 

\subsection{Biological Significance of Blind Rediscovery}
The blind rediscovery of Palbociclib in breast cancer, Erlotinib in non-small cell lung cancer, and Gefitinib in glioblastoma provides empirical validation of the underlying hypothesis: when malignant cells undergo oncogenic transformation, their global transcriptional output represents a coherent physiological state. Small molecules that invert this signature do not merely act as broad poisons; rather, their induced transcriptional cascades directly counteract the hallmark transcriptomic modules necessary for tumor proliferation and survival.

\subsection{Novel Hypothesis Generation for Refractory Tumors}
Beyond positive-control validation, our comparative results yield testable preclinical hypotheses for disease types with limited therapeutic options:
\begin{enumerate}
    \item \textbf{CDK4/6 Inhibition in Squamous Cell Carcinomas}: The high reversal score of Palbociclib in both Lung Squamous Cell ($-0.364$) and Head and Neck Squamous Cell ($-0.039$) supports emerging clinical trial data evaluating dual CDK4/6 and immune checkpoint inhibition in squamous histologies.
    \item \textbf{MEK/MAPK Blockade in Prostate Cancer}: Trametinib's prioritization in TCGA-PRAD suggests that subsets of prostate adenocarcinomas characterized by high neuroendocrine or MAPK pathway activity may benefit from targeted MEK inhibition.
    \item \textbf{RNA Polymerase I Inhibition in Liver Cancer}: The top hit in TCGA-LIHC, \textbf{CX-5461}, is a selective inhibitor of RNA polymerase I transcription, highlighting ribosomal biogenesis as a vulnerability in hepatocellular carcinoma.
\end{enumerate}

\subsection{Methodological Limitations}
While powerful, several limitations must be considered:
\begin{enumerate}
    \item \textbf{Cellular Composition and Bulk Deconvolution}: TCGA specimens are bulk tissue homogenates containing stromal and immune infiltrating populations. Some observed differential expression reflects cellular proportion shifts rather than pure cell-autonomous reprogramming.
    \item \textbf{LINCS L1000 Landmark Space}: LINCS directly profiles 978 landmark genes, inferring remaining transcripts via linear regression. While directional sets capture major pathways, splice isoforms and non-coding RNAs are unrepresented.
    \item \textbf{In Vitro Translation}: In vitro drug responses in cell culture models do not account for in vivo pharmacokinetics, bioavailability, tumor microenvironment penetration, or host toxicity.
\end{enumerate}

\section{Conclusion and Future Work}
We have demonstrated a mathematically rigorous, fully automated, and reproducible framework for computational drug reversal across ten major cancer types. By strictly enforcing matched non-malignant controls and permutation significance, IRIS successfully rediscovers frontline targeted antineoplastics and generates robust preclinical hypotheses for therapeutic repurposing. Future extensions will incorporate single-cell RNA sequencing deconvolution and cross-study normalization against GTEx tissue compendia to expand comparative reversal to hematologic malignancies.

\section*{Data and Code Availability}
The IRIS computational framework, raw analysis pipelines, Next.js interactive visualization platform, and all precomputed genomic artifacts are open source and publicly accessible at:
\begin{center}
\url{https://github.com/NotNxel/IRIS-Althea}
\end{center}

\begin{thebibliography}{00}
\bibitem{dimasi2016innovation} J.~A. DiMasi, H.~G. Grabowski, and R.~W. Hansen, ``Innovation in the pharmaceutical industry: New estimates of R\&D costs,'' \emph{Journal of Health Economics}, vol.~47, pp.~20--33, 2016.
\bibitem{vogelstein2013cancer} B.~Vogelstein et al., ``Cancer genome landscapes,'' \emph{Science}, vol.~339, no.~6127, pp.~1546--1558, 2013.
\bibitem{lamb2006connectivity} J.~Lamb et al., ``The Connectivity Map: Using gene-expression signatures to connect small molecules, genes, and disease,'' \emph{Science}, vol.~313, no.~5795, pp.~1929--1935, 2006.
\bibitem{subramanian2017lincs} A.~Subramanian et al., ``A Next Generation Connectivity Map: L1000 Platform and the First 1,000,000 Profiles,'' \emph{Cell}, vol.~171, no.~6, pp.~1437--1452, 2017.
\bibitem{dobin2013star} A.~Dobin et al., ``STAR: ultrafast universal RNA-seq aligner,'' \emph{Bioinformatics}, vol.~29, no.~1, pp.~15--21, 2013.
\bibitem{benjamini1995controlling} Y.~Benjamini and Y.~Hochberg, ``Controlling the false discovery rate: a practical and powerful approach to multiple testing,'' \emph{Journal of the Royal Statistical Society: Series B}, vol.~57, no.~1, pp.~289--300, 1995.
\bibitem{weinstein2013cancer} J.~N. Weinstein et al., ``The Cancer Genome Atlas Pan-Cancer analysis project,'' \emph{Nature Genetics}, vol.~45, no.~10, pp.~1113--1120, 2013.
\bibitem{chen2013enrichr} E.~Y. Chen et al., ``Enrichr: interactive and collaborative HTML5 gene list enrichment analysis tool,'' \emph{BMC Bioinformatics}, vol.~14, no.~1, p.~128, 2013.
\bibitem{welch1947generalization} B.~L. Welch, ``The generalization of `Student's' problem when several different population variances are involved,'' \emph{Biometrika}, vol.~34, no.~1/2, pp.~28--35, 1947.
\bibitem{gtex2020genomic} GTEx Consortium, ``The GTEx Consortium atlas of genetic regulatory effects across human tissues,'' \emph{Science}, vol.~369, no.~6509, pp.~1318--1330, 2020.
\end{thebibliography}

\end{document}
